
\documentclass[tikz,border=5pt]{standalone}
\usepackage[french]{babel}
\usepackage{amsmath}
\usepackage{amssymb}
\usepackage{xcolor}
\usetikzlibrary{calc}
\usepackage{pgfplots}
\pgfplotsset{compat=1.17}

\usepackage{tkz-euclide} % Use tikz-euclide for geometry
\begin{document}

% Exercise Text
\section*{139 \quad \textcolor{teal}{Théorème de Pappus} \hfill \textcolor{green!50!black}{Histoire des sciences}}

On considère les droites d'équations réduites $y = x$ et $y = -3$. On place sur la première droite les points $A$, $B$ et $C$ d'abscisses respectives $0$, $1$ et $4$. Puis sur la seconde droite on place les points $D$, $E$ et $F$ d'abscisses respectives $1$, $4$ et $7$.

\begin{enumerate}
    \item Déterminer les équations cartésiennes des droites $(BF)$ et $(CE)$.
    \item Déterminer les coordonnées du point d'intersection $M$ des droites $(BF)$ et $(CE)$.
    \item Déterminer les équations cartésiennes des droites $(AF)$ et $(CD)$.
    \item Déterminer les coordonnées du point d'intersection $N$ des droites $(AF)$ et $(CD)$.
    \item Déterminer les équations cartésiennes des droites $(AE)$ et $(BD)$.
    \item Déterminer les coordonnées du point d'intersection $P$ des droites $(AE)$ et $(BD)$.
    \item \textbf{Démontrer que les points $M$, $N$ et $P$ sont alignés.}
\end{enumerate}

\textcolor{green!50!black}{\textbf{Remarque}} Cette propriété est vraie quelle que soit la position des points $A$, $B$ et $C$ sur une droite et des points $D$, $E$ et $F$ sur une autre droite.






\begin{tikzpicture}
% Define the axis
\begin{axis}[
    width=10cm, height=8cm,
    axis lines=middle,
    xmin=-1, xmax=8,
    ymin=-4, ymax=5,
    grid=major,
    xtick={0,1,2,3,4,5,6,7},
    ytick={-3,-2,-1,0,1,2,3,4,5},
    xlabel={$x$}, ylabel={$y$},
    axis equal,
    tick label style={font=\scriptsize},
    label style={font=\scriptsize}
]

% Define points
\coordinate (A) at (0,0);
\coordinate (B) at (1,1);
\coordinate (C) at (4,4);
\coordinate (D) at (1,-3);
\coordinate (E) at (4,-3);
\coordinate (F) at (7,-3);

% Plot the points
\addplot[only marks, mark=*] coordinates {(0,0) (1,1) (4,4) (1,-3) (4,-3) (7,-3)};

% Add labels
\node[below left] at (A) {$A$};
\node[above right] at (B) {$B$};
\node[above right] at (C) {$C$};
\node[below left] at (D) {$D$};
\node[below right] at (E) {$E$};
\node[below right] at (F) {$F$};

% Draw the main lines
\addplot[blue] coordinates {(0,0) (8,8)}; % Line y = x
\addplot[red] coordinates {(-1,-3) (8,-3)}; % Line y = -3

% BF and CE lines
\draw[black, thick] (B) -- (F);
\draw[green!50!black, thick] (C) -- (E);

% AF and CD lines
\draw[blue, thick] (A) -- (F);
\draw[red, thick] (C) -- (D);

% AE and BD lines
\draw[orange, thick] (A) -- (E);
\draw[magenta, thick] (B) -- (D);

% Calculate M (intersection of BF and CE)
\tkzInterLL(B,F)(C,E) \tkzGetPoint{M}
\node[fill=black, circle, inner sep=1.5pt] at (M) {};
\node[above right] at (M) {$M$};

% Calculate N (intersection of AF and CD)
\tkzInterLL(A,F)(C,D) \tkzGetPoint{N}
\node[fill=black, circle, inner sep=1.5pt] at (N) {};
\node[below left] at (N) {$N$};

% Calculate P (intersection of AE and BD)
\tkzInterLL(A,E)(B,D) \tkzGetPoint{P}
\node[fill=black, circle, inner sep=1.5pt] at (P) {};
\node[below left] at (P) {$P$};

\end{axis}
\end{tikzpicture}

\end{document}
